%% Complete proofs under ZF plus the Boolean Prime Ideal Theorem.

\section{Stone representation and duality}

\begin{annexitem}{annex:p10.stone.1}{thm:p10-ultrafilter-separation}{Ultrafilter dichotomy and separation}
\begin{proof}
For a filter \(F\) and element \(a\), the filter generated by
\(F\cup\{a\}\) is
\[
 F[a]=\{b:\text{some }f\in F\text{ satisfies }f\wedge a\leq b\}.
\]
It contains \(F\) and \(a\), contains \(1\), and is upward closed.
If \(f\wedge a\leq b\) and \(g\wedge a\leq c\), then
\((f\wedge g)\wedge a\leq b\wedge c\), proving meet closure.
Any filter containing \(F\) and \(a\) contains each of these elements,
so this is the generated filter.

If \(U\) is maximal proper and \(a\notin U\), then \(U[a]\) is
improper. Thus some \(u\in U\) has \(u\wedge a=0\).
In a Boolean algebra this is equivalent to \(u\leq\neg a\), so
\(\neg a\in U\).
Both \(a,\neg a\) cannot belong to a proper filter, because their
meet is \(0\). This proves the dichotomy.
Conversely, suppose a proper filter has the dichotomy.
Any strictly larger filter contains an \(a\notin U\) as well as
\(\neg a\in U\), so it is improper. Therefore \(U\) is maximal proper.

If either \(a\) or \(b\) belongs to \(U\), upward closure puts
\(a\vee b\) in \(U\).
If neither belongs, the dichotomy puts both complements in \(U\),
and hence puts \(\neg(a\vee b)\) in \(U\); properness excludes
\(a\vee b\). This proves primeness, and induction gives the
corresponding assertion for any finite join, including the empty join.

Now let \(b\notin F\). The filter \(F[\neg b]\) is proper:
otherwise some \(f\in F\) would satisfy \(f\wedge\neg b=0\),
so \(f\leq b\), contradicting \(b\notin F\).
BPIT extends this proper filter to an ultrafilter \(U\).
It contains \(\neg b\), hence omits \(b\), and contains \(F\).
For a nonzero \(a\), the principal filter
\(\{b:a\leq b\}\) is proper, so BPIT yields an ultrafilter containing
\(a\). If \(a\ne b\), at least one of
\(a\wedge\neg b\) and \(b\wedge\neg a\) is nonzero.
An ultrafilter containing that element contains exactly one of \(a,b\).
Finally, when \(0=1\), every filter contains \(0\), so no proper
filter or ultrafilter exists.
\end{proof}
\end{annexitem}

\begin{annexitem}{annex:p10.stone.2}{thm:p10-stone-representation}{Stone representation theorem}
\begin{proof}
\paragraph{The basic clopens and separation.}
The ultrafilter laws imply
\[
 \widehat0=\varnothing,\quad \widehat1=X_B,\quad
 \widehat{a\wedge b}=\widehat a\cap\widehat b,\quad
 \widehat{a\vee b}=\widehat a\cup\widehat b,\quad
 \widehat{\neg a}=X_B\setminus\widehat a.
\]
Thus the displayed basis consists of clopens.
Distinct ultrafilters \(U,V\) disagree on some \(a\).
Interchanging them if needed, take \(a\in U\setminus V\).
Then \(\widehat a\) and \(\widehat{\neg a}\) are disjoint open
neighborhoods of \(U,V\), proving the Hausdorff property.

\paragraph{Compactness for arbitrary open covers.}
Let \(\mathcal U\) be an open cover of \(X_B\).
Form the set of all subordinate basic clopens, represented by
\[
 D=\{a\in B:\widehat a\subseteq O
                      \text{ for some }O\in\mathcal U\}.
\]
The family \(\{\widehat a:a\in D\}\) covers \(X_B\), because the
basic clopens form a basis. No selection of one neighborhood per
point is involved.
Suppose \(\mathcal U\) has no finite subcover.
For every finite \(E\subseteq D\), we must have
\(\bigvee E\ne1\).
Otherwise its basic clopens would cover \(X_B\); choosing a
containing member of \(\mathcal U\) for each of the finitely many
members of \(E\) would give a finite subcover.

Consequently each finite meet of elements of
\(\{\neg a:a\in D\}\) is nonzero.
The filter they generate is
\[
 F=\left\{b:\text{some finite }E\subseteq D
                  \text{ satisfies }\bigwedge_{a\in E}\neg a\leq b\right\}.
\]
The empty meet is \(1\); the displayed nonzero property shows
\(0\notin F\).
Union of two finite index sets proves meet closure, and upward
closure is immediate.
Thus \(F\) is proper.
BPIT supplies an ultrafilter \(U\supseteq F\).
It contains \(\neg a\) for every \(a\in D\), so it belongs to
none of the basic clopens covering \(X_B\), a contradiction.
Every open cover therefore has a finite subcover.
Together with the preceding paragraph, \(X_B\) is a Stone space.

\paragraph{The algebra representation is onto.}
The basic identities make \(\beta_B\) a Boolean homomorphism.
Ultrafilter separation shows it is injective.
Let \(K\) be any clopen subset of \(X_B\).
The family of all basic clopens contained in \(K\), together with
the open set \(X_B\setminus K\), covers \(X_B\).
Compactness gives finitely many basic clopens
\(\widehat{a_1},\ldots,\widehat{a_n}\subseteq K\) covering \(K\).
Then
\[
 K=\widehat{a_1}\cup\cdots\cup\widehat{a_n}
   =\widehat{a_1\vee\cdots\vee a_n}.
\]
For \(K=\varnothing\), the empty join \(0\) suffices.
Thus every clopen is in the image.

When \(B\) is trivial, \(X_B=\varnothing\), whose topology is compact,
Hausdorff, and has a clopen basis. Both \(B\) and
\(\operatorname{Clop}(\varnothing)\) have one element, and the
same map is an isomorphism. This also verifies that edge case directly.
\end{proof}
\end{annexitem}

\begin{annexitem}{annex:p10.stone.3}{thm:p10-stone-space-reconstruction}{Recovery from the clopen algebra}
\begin{proof}
For \(x\in X\), its clopen neighborhoods contain \(X\), exclude
\(\varnothing\), and are upward closed and closed under meets.
Exactly one of a clopen and its complement contains \(x\).
The ultrafilter criterion therefore proves that \(\varepsilon_X(x)\)
is an ultrafilter.

Let \(U\) be any ultrafilter of \(\operatorname{Clop}(X)\).
Its clopens have the finite-intersection property: every finite
intersection belongs to \(U\), and none is empty.
Compactness implies their total intersection is nonempty.
Indeed, if that intersection were empty, their open complements
would cover \(X\), and a finite subcover would give an empty finite
intersection.
Choose \(x\) in the total intersection.
Every \(C\in U\) contains \(x\).
Conversely, if a clopen \(C\) contains \(x\) and is absent from \(U\),
its complement lies in \(U\), contradicting the choice of \(x\).
Hence \(U=\varepsilon_X(x)\).

Clopens separate distinct points: the Hausdorff property gives
disjoint open neighborhoods, and a basic clopen inside either one
contains its selected point and excludes the other.
Thus two points cannot yield the same ultrafilter, proving
bijectivity.
For every \(C\in\operatorname{Clop}(X)\),
\[
 \varepsilon_X^{-1}[\widehat C]=C,
 \qquad \varepsilon_X[C]=\widehat C.
\]
The first equality proves continuity on a basis; the second proves
that the bijection is open, again using the clopen basis.
Its inverse is continuous, so it is a homeomorphism.
If \(X=\varnothing\), its clopen algebra is trivial and has no
ultrafilters. The unique empty map is the required homeomorphism.
\end{proof}
\end{annexitem}

\begin{annexitem}{annex:p10.stone.4}{thm:p10-stone-duality}{Functorial Stone duality}
\begin{proof}
\paragraph{Maps are well defined.}
For a Boolean homomorphism \(h:B\to C\) and \(U\in X_C\),
the inverse image \(h^{-1}[U]\) contains \(1\), excludes \(0\),
and is closed upward and under meets.
Preservation of complement and the ultrafilter criterion show
it selects exactly one of \(a,\neg a\) for each \(a\in B\).
Thus it belongs to \(X_B\).
For a basic clopen,
\[
 X(h)^{-1}[\widehat a]=\widehat{h(a)}.
\]
This proves continuity.
Conversely, for a continuous \(g:X\to Y\), the inverse image of a
clopen is clopen. The usual inverse-image identities preserve
intersection, union, complement, and both bounds, so
\(\operatorname{Clop}(g)\) is a Boolean homomorphism.

Inverse images preserve identity maps and satisfy
\((k\circ h)^{-1}=h^{-1}\circ k^{-1}\).
Applying this equality to the appropriate subsets gives the
identity and reversed-composition laws for both assignments.
The object assignments take values in the required categories by
the two preceding theorems.

\paragraph{Naturality.}
For \(a\in B\), the preceding basic-preimage identity gives
\[
 \operatorname{Clop}(X(h))(\beta_B(a))
     =\widehat{h(a)}=\beta_C(h(a)).
\]
For \(x\in X\) and \(D\in\operatorname{Clop}(Y)\),
\[
\begin{aligned}
 D\in X(\operatorname{Clop}(g))(\varepsilon_X(x))
 &\Longleftrightarrow g^{-1}[D]\in\varepsilon_X(x)\\
 &\Longleftrightarrow x\in g^{-1}[D]\\
 &\Longleftrightarrow g(x)\in D\\
 &\Longleftrightarrow D\in\varepsilon_Y(g(x)).
\end{aligned}
\]
These calculations prove both naturality equations.

\paragraph{The two inverse identities.}
For \(U\in X_B\) and \(a\in B\),
\[
 a\in X(\beta_B)(\varepsilon_{X_B}(U))
 \ \Longleftrightarrow\ 
 \widehat a\in\varepsilon_{X_B}(U)
 \ \Longleftrightarrow\ U\in\widehat a
 \ \Longleftrightarrow\ a\in U.
\]
For \(D\in\operatorname{Clop}(X)\),
\[
 \operatorname{Clop}(\varepsilon_X)
       (\beta_{\operatorname{Clop}(X)}(D))
   =\{x:D\in\varepsilon_X(x)\}=D.
\]
Thus both composites are identities.
Since every \(\beta_B\) is an isomorphism and every
\(\varepsilon_X\) is a homeomorphism, the natural reconstruction
maps give the claimed contravariant equivalence.

\paragraph{Recovering a homomorphism from a continuous map.}
Given continuous \(f:X_C\to X_B\), put
\[
 h=\beta_C^{-1}\circ\operatorname{Clop}(f)\circ\beta_B.
\]
It is a Boolean homomorphism.
For \(U\in X_C\) and \(a\in B\),
\[
 a\in X(h)(U)
 \ \Longleftrightarrow\ h(a)\in U
 \ \Longleftrightarrow\ U\in f^{-1}[\widehat a]
 \ \Longleftrightarrow\ a\in f(U).
\]
Hence \(X(h)=f\).
If another homomorphism has that same dual map, the basic-preimage
identity gives the same clopen \(\widehat{h(a)}\) for each \(a\).
Injectivity of \(\beta_C\) makes its values equal to those of \(h\),
proving uniqueness.
All constructions cover empty Stone spaces and trivial Boolean
algebras, with the same inverse-image identities.
\end{proof}
\end{annexitem}
